% ============================================================
%  body.tex - Experiment B (pi-0.5 on LIBERO)
%  Frames only (no preamble).  Imported by
%      slides/2026.10.05-2026.10.11-pi05-libero/pi05-libero-slides.tex
%      slides.tex
%  The importing document must load the `ego-beamer` style first.
% ============================================================

\egodeckid{2026.10.05-2026.10.11-pi05-libero}

\title{Reproducing \texorpdfstring{$\pi_{0.5}$}{pi-0.5} on LIBERO}
\subtitle{Official OpenPI weights and protocol, single A100}
\author{Yutao-Yang}
\date{2026.10.05--2026.10.11}

\egotitleframe%
\egooutline%

% ============================================================
\section{Motivation}
% ============================================================

\begin{frame}{Why reproduce \texorpdfstring{$\pi_{0.5}$}{pi-0.5}?}
  \begin{itemize}
    \item \texorpdfstring{$\pi_{0.5}$}{pi-0.5} is a vision--language--action
          model: two camera streams and a language task in, a continuous
          $10 \times 7$ action chunk out.
    \item The authors report \textbf{96.85\,\%} over the four standard LIBERO
          suites.
    \item Question: does the \emph{released} checkpoint, driven by the
          \emph{official} OpenPI implementation, reproduce that number on one
          A100?
  \end{itemize}
  \vfill
  \begin{block}{What is held fixed}
    Released LIBERO weights, authors' task protocol, environment and budgets.
    No retraining, no per-seed selection. A P10 highest-precision run is kept
    as a \emph{control}, not as the main table.
  \end{block}
\end{frame}

% ============================================================
\section{Benchmark and setup}
% ============================================================

\begin{frame}{The four LIBERO suites}
  \begin{columns}[T]
    \begin{column}{0.46\linewidth}
      \begin{itemize}
        \item $40$ tasks $\times$ $50$ fixed initial states $=$
              \textbf{2000} closed-loop episodes.
        \item 20\,Hz control; $256$ RGB frames, $180^{\circ}$ rotation, padded
              to $224$.
        \item Success is the official \texttt{done}/\texttt{goal} signal;
              timeouts count as failures.
      \end{itemize}
    \end{column}
    \begin{column}{0.5\linewidth}
      \centering
      \vspace{-2mm}
      \begin{tabular}{@{}llr@{}}
        \toprule
        Suite & Focus & Official \\
        \midrule
        Spatial   & spatial relations   & 98.8\,\% \\
        Object    & target objects      & 98.2\,\% \\
        Goal      & goal states         & 98.0\,\% \\
        LIBERO-10 & longer horizon      & 92.4\,\% \\
        \midrule
        Mean      & 40 tasks               & 96.85\,\% \\
        \bottomrule
      \end{tabular}
      \vfill
      {\scriptsize Official numbers come from the OpenPI release table, not
       from this run.}
    \end{column}
  \end{columns}
\end{frame}

% ============================================================
\section{Result}
% ============================================================

\begin{frame}{Headline result}
  \begin{columns}[T]
    \begin{column}{0.5\linewidth}
      \begin{itemize}
        \item Reproduced: $\mathbf{1935/2000 = 96.75\,\%}$.
        \item Official: $96.85\,\%$ ($-0.10$ percentage points).
        \item All 2000 episodes retained; all $64\,761$ inferences and
              $320\,142$ executed steps finite.
      \end{itemize}
    \end{column}
    \begin{column}{0.46\linewidth}
      \centering
      \includegraphics[width=\linewidth]{2026.10.05-2026.10.11-pi05-libero/official-control-native.png}
    \end{column}
  \end{columns}
  \vfill
  \begin{block}{Per suite (success / 500)}
    Spatial $495$ \quad Object $490$ \quad Goal $491$ \quad LIBERO-10 $459$.
    \hfill {\normalfont{}P10 precision control: $97.00\,\%$.}
  \end{block}
\end{frame}

\begin{frame}{Every task}
  \centering
  \includegraphics[width=0.9\linewidth]{2026.10.05-2026.10.11-pi05-libero/all-40-task-rates.png}
\end{frame}

% ============================================================
\section{Compute}
% ============================================================

\begin{frame}{Compute and control budget}
  \begin{columns}[T]
    \begin{column}{0.44\linewidth}
      \begin{itemize}
        \item End-to-end wall time: $9673.44$\,s ($\approx 2$\,h\,41\,m) ---
              the four suites overlap, so they must not be summed.
        \item $64\,761$ model queries; $320\,142$ executed control steps.
      \end{itemize}
    \end{column}
    \begin{column}{0.52\linewidth}
      \centering
      \vspace{-1mm}
      \begin{tabular}{@{}lrr@{}}
        \toprule
        Suite & Wall & Steps \\
        \midrule
        spatial   & 85.36 min  & 52\,625  \\
        object    & 95.89 min  & 70\,938  \\
        goal      & 85.43 min  & 57\,342  \\
        LIBERO-10 & 147.45 min & 139\,237 \\
        \bottomrule
      \end{tabular}
      \vfill
      {\scriptsize Predict 10 actions, execute the first 5, then replan.}
    \end{column}
  \end{columns}
\end{frame}

% ============================================================
\section{Teaching cases}
% ============================================================

\begin{frame}{Genuine model inputs: external and wrist views}
  \centering
  \begin{tabular}{cccc}
    \includegraphics[width=0.18\linewidth]{2026.10.05-2026.10.11-pi05-libero/spatial-t0-i0-first-agent.png} &
    \includegraphics[width=0.18\linewidth]{2026.10.05-2026.10.11-pi05-libero/object-t0-i0-first-agent.png} &
    \includegraphics[width=0.18\linewidth]{2026.10.05-2026.10.11-pi05-libero/goal-t0-i0-first-agent.png} &
    \includegraphics[width=0.18\linewidth]{2026.10.05-2026.10.11-pi05-libero/goal-t0-i1-first-agent.png} \\[1mm]
    \includegraphics[width=0.18\linewidth]{2026.10.05-2026.10.11-pi05-libero/spatial-t0-i0-first-wrist.png} &
    \includegraphics[width=0.18\linewidth]{2026.10.05-2026.10.11-pi05-libero/object-t0-i0-first-wrist.png} &
    \includegraphics[width=0.18\linewidth]{2026.10.05-2026.10.11-pi05-libero/goal-t0-i0-first-wrist.png} &
    \includegraphics[width=0.18\linewidth]{2026.10.05-2026.10.11-pi05-libero/goal-t0-i1-first-wrist.png} \\
    {\scriptsize spatial} & {\scriptsize object} & {\scriptsize goal i0} & {\scriptsize goal i1}
  \end{tabular}
  \vfill
  {\scriptsize First real query of four teaching cases. These six P9 cases are
   engineering cases and do not replace the 2000-episode score.}
\end{frame}

\begin{frame}{Executed actions}
  \vspace{-1mm}
  \centering
  \begin{tabular}{cc}
    \includegraphics[width=0.30\linewidth]{2026.10.05-2026.10.11-pi05-libero/spatial-t0-i0-actions.png} &
    \includegraphics[width=0.30\linewidth]{2026.10.05-2026.10.11-pi05-libero/object-t0-i0-actions.png} \\[0.5mm]
    {\scriptsize spatial t0 i0 (success, 101 steps)} &
    {\scriptsize object t0 i0 (success, 131 steps)} \\[1.5mm]
    \includegraphics[width=0.30\linewidth]{2026.10.05-2026.10.11-pi05-libero/goal-t0-i0-actions.png} &
    \includegraphics[width=0.30\linewidth]{2026.10.05-2026.10.11-pi05-libero/goal-t0-i1-actions.png} \\[0.5mm]
    {\scriptsize goal t0 i0 (success, 116 steps)} &
    {\scriptsize goal t0 i1 (timeout failure, 300 steps)}
  \end{tabular}
\end{frame}

\begin{frame}{Rollouts}
  \begin{columns}[T]
    \begin{column}{0.48\linewidth}
      \centering
      \includegraphics[width=0.62\linewidth]{2026.10.05-2026.10.11-pi05-libero/spatial-t0-i0-last-agent.png}
      \vfill
      \egovideolink{2026.10.05-2026.10.11-pi05-libero/spatial-t0-i0-rollout.mp4}\\[1mm]
      {\scriptsize spatial t0 i0 --- success}
    \end{column}
    \begin{column}{0.48\linewidth}
      \centering
      \includegraphics[width=0.62\linewidth]{2026.10.05-2026.10.11-pi05-libero/goal-t0-i1-last-agent.png}
      \vfill
      \egovideolink{2026.10.05-2026.10.11-pi05-libero/goal-t0-i1-rollout.mp4}\\[1mm]
      {\scriptsize goal t0 i1 --- timeout failure}
    \end{column}
  \end{columns}
  \vfill
  {\scriptsize Last query of each run. Videos are written at 10\,fps for viewing,
   so they last about twice the simulated action time. Click a button to open
   the file next to the PDF.}
\end{frame}

% ============================================================
\section{Normal failures}
% ============================================================

\begin{frame}{An open drawer is not success}
  \begin{columns}[T]
    \begin{column}{0.27\linewidth}
      \centering
      \includegraphics[width=\linewidth]{2026.10.05-2026.10.11-pi05-libero/drawer-success-agent.png}\\[1mm]
      {\scriptsize ``open the middle drawer'' --- \textbf{success}\\
       middle drawer $-0.14426769$}
    \end{column}
    \hfill
    \begin{column}{0.27\linewidth}
      \centering
      \includegraphics[width=\linewidth]{2026.10.05-2026.10.11-pi05-libero/drawer-failure-agent.png}\\[1mm]
      {\scriptsize ``open the middle drawer'' --- \textbf{failure}\\
       middle drawer $0.00147690$\\
       bottom drawer opened}
    \end{column}
  \end{columns}
  \vfill
  {\scriptsize The BDDL goal requires the \emph{middle}-drawer joint below
   $-0.14$. Both rolls were replayed on CPU with a maximum state error of $0$
   and agreed \texttt{done}/\texttt{goal} signals.}
\end{frame}

% ============================================================
\section{Takeaway}
% ============================================================

\begin{frame}{What this shows --- and what it does not}
  \begin{itemize}
    \item The released checkpoint \textbf{reproduces the official score} to
          within $0.10$ percentage points using the official implementation.
    \item Agreement is about the aggregate number: the release publishes no
          per-episode predictions, so per-sample identity is not claimed.
    \item The P12/P10 gap is a single-run precision control and is \emph{not}
          a stable precision advantage.
    \item Non-finite engineering failures (NaN GPU inputs) are reported
          separately from genuine task failures such as the drawer case.
    \item In this configuration the 8-D state is \emph{not} a network
          condition, so no state-conditioned claim is made.
    \item This is a baseline reproduction, not a fair architecture comparison
          with the other model in this project.
  \end{itemize}
\end{frame}

\begin{frame}{Reference}
  \begin{itemize}
    \item \texorpdfstring{$\pi_{0.5}$}{pi-0.5}: Black et al., \emph{A
          Vision--Language--Action Model with Open-World Generalization},
          arXiv:2504.16054.
    \item OpenPI reference implementation: revision
          \texttt{215abfb2}.
    \item LIBERO: Liu et al., \emph{Benchmarking Knowledge Transfer for
          Lifelong Robot Learning}, arXiv:2306.03310.
  \end{itemize}
  \vfill
  \centering
  {\small Full report: \texttt{reports/2026.10.05-2026.10.11-pi05-libero/}}
\end{frame}
